% Generated from 2/2-8.md; technical content preserved.
\exercise{۲.۸}
تمرین بعدی از لیست تکالیف، \textbf{\lr{Exercise 2.8}} از فصل دوم کتاب مرجع است.

\begin{quote}
\textbf{صورت مسئله (\lr{Exercise 2.8}):}

آیا اثر بدنه (\textbf{\lr{Body Effect}}) در یک فرآیند ساخت، تعداد ترانزیستورهایی را که می‌توان در فرکانس‌های پایین به‌صورت سری در یک گیت \lr{CMOS} قرار داد، محدود می‌کند؟

\end{quote}

\subsection*{۱. پاسخ مستقیم (\lr{Direct Answer})}
\textbf{خیر (\lr{No}).} اثر بدنه در فرکانس‌های بسیار پایین یا عملکرد \lr{DC}، مانع فیزیکی یا منطقی برای قرارگیری تعداد نامحدود ترانزیستور به‌صورت سری ایجاد نمی‌کند.

\subsection*{۲. تحلیل فیزیکی و مداری در فرکانس پایین (\lr{DC} / \lr{Low-Frequency Analysis})}

برای درک رفتار استک سری (\lr{Series Stack})، شبکه \lr{Pull-Down} یک گیت با \(N\) ترانزیستور سری \lr{nMOS} را در نظر بگیرید که خروجی را به زمین متصل می‌کند:

\begin{enumerate}
\item \textbf{شرایط حالت ماندگار (\lr{Steady-State Operation}):}
در فرکانس‌های پایین، مدار زمان کافی برای شارژ و دشارژ کامل خازن‌های پارازیتی گره‌ها دارد. وقتی تمام گیت‌های این ترانزیستورها به ولتاژ \(V_{DD}\) متصل باشند، هدایت الکتریکی در کل شاخه برقرار می‌شود.

\item \textbf{تخلیه گره‌های میانی (\lr{Discharge of Internal Nodes}):}
خروجی گیت از طریق استک ترانزیستورها تخلیه شده و پتانسیل آن در نهایت به سمت صفر ولت (\(0\text{\lr{ V}}\)) میل می‌کند. به همین ترتیب، جریان نشتی یا تخلیه خازنی باعث می‌شود که پتانسیل تمام گره‌های میانی بین ترانزیستورها (\(V_{internal}\)) نیز در حالت ماندگار به صفر ولت برسد:
\[V_{S,i} \to 0\text{\lr{ V}} \quad (\text{\lr{for all }} i = 1, 2, \dots, N)\]

\item \textbf{حذف ولتاژ سورس-بدنه در حالت تعادل:}
چون سابستریت (بدنه ترانزیستورهای \lr{nMOS}) به زمین متصل است (\(V_B = 0\text{\lr{ V}}\))، در حالت ماندگار ولتاژ سورس-بدنه تمام ترانزیستورها برابر صفر خواهد شد:
\[V_{SB,i} = V_{S,i} - V_B = 0 - 0 = 0\text{\lr{ V}}\]

\item \textbf{مقدار ولتاژ آستانه در پایان گذار:}
طبق رابطه اثر بدنه:
\[V_t = V_{t0} + \gamma \left( \sqrt{\phi_s + V_{SB}} - \sqrt{\phi_s} \right)\]
با رسیدن \(V_{SB}\) به صفر، مقدار \(V_t\) تمام ترانزیستورهای شاخه به ولتاژ آستانه نامی اولیه (\(V_{t0}\)) بازمی‌گردد. بنابراین هیچ ترانزیستوری در حالت ماندگار به دلیل افزایش \(V_t\) خاموش نمی‌شود و مسیر هدایت \lr{DC} برقرار مانده و سطح منطقی کامل صفر (\lr{Strong 0}) روی گره خروجی تحویل داده می‌شود.

\end{enumerate}

\subsection*{۳. عامل محدودکننده واقعی چیست؟ (\lr{Practical Limitations})}

اگرچه در فرکانس پایین محدودیت عملکرد منطقی وجود ندارد، اما در فرکانس‌های عملیاتی و کاربردهای واقعی، قرار دادن ترانزیستورهای سری با محدودیت‌های شدید زیر روبه‌رو است:

\begin{enumerate}
\item \textbf{رشد درجه دوم تأخیر انتشار (\lr{Quadratic Delay Scaling}):}
طبق مدل تأخیر \lr{Elmore}، با افزایش تعداد ترانزیستورهای سری (\(N\))، علاوه بر مقاومت مسیر، ظرفیت خازنی گره‌های میانی نیز افزایش می‌یابد؛ در نتیجه تأخیر انتشار متناسب با توان دوم تعداد ترانزیستورها رشد می‌کند:
\[t_{pd} \propto N^2\]

\item \textbf{تأثیر دینامیکی اثر بدنه در طول گذار (\lr{Transient Degradation}):}
در حین سوئیچینگ، تا پیش از تخلیه کامل بار گره‌های داخلی، پتانسیل سورس ترانزیستورهای بالایی مثبت است (\(V_{SB} > 0\)). این امر به‌طور موقت ولتاژ آستانه آن‌ها را افزایش داده و جریان درین را تضعیف می‌کند؛ در نتیجه سرعت گذار خروجی بیش از پیش کاهش می‌یابد.

\item \textbf{توصیه صنعتی در طراحی \lr{VLSI}:}
به دلیل افت شدید کارایی و افزایش ابعاد فیزیکی برای جبران مقاومت ترانزیستورها، در کتاب مرجع تأکید شده است که عملاً نباید بیش از ۴ (یا حداکثر ۵) ترانزیستور را در یک گیت به‌صورت سری قرار داد و در صورت نیاز به تعداد ورودی بیشتر، از ساختارهای درختی گیت‌های کوچک‌تر (مانند ساختار درختی \lr{NAND}) استفاده می‌شود.

\end{enumerate}
